% Auto-generated Academic Glossary with Bidirectional Hyperlinks
\begin{longtable}{p{2.8cm} p{4.2cm} p{3.6cm}}
\toprule
\textbf{English Term / 学术用语} & \textbf{Chinese Definition / 中文规范释义} & \textbf{Academic Context / 语境注记} \\
\midrule
\endhead
\bottomrule
\endfoot
\hypertarget{glossary:almost disjoint}{}\hyperlink{termfirst:almost disjoint}{\textbf{almost disjoint}} & {\kaishu 几乎不相交} & {\small\textcolor{black!70}{交集测度为零的几何族}} \\
\hypertarget{glossary:by virtue of}{}\hyperlink{termfirst:by virtue of}{\textbf{by virtue of}} & {\kaishu 依据 / 凭借} & {\small\textcolor{black!70}{高频数学论证学术表达}} \\
\hypertarget{glossary:exterior measure}{}\hyperlink{termfirst:exterior measure}{\textbf{exterior measure}} & {\kaishu 外测度} & {\small\textcolor{black!70}{全空间子集上的广义外测度}} \\
\hypertarget{glossary:measurable set}{}\hyperlink{termfirst:measurable set}{\textbf{measurable set}} & {\kaishu 可测集} & {\small\textcolor{black!70}{依据外测度能赋予测度的集合}} \\
\hypertarget{glossary:pointwise limit}{}\hyperlink{termfirst:pointwise limit}{\textbf{pointwise limit}} & {\kaishu 逐点极限} & {\small\textcolor{black!70}{函数序列在每个点处的极限}} \\
\end{longtable}
